\documentclass[a4paper,twoside]{article}

\usepackage{morewrites}

\usepackage[svgnames,dvipsnames,x11names]{xcolor}
\usepackage{graphicx}
\usepackage[english]{babel}
\usepackage[colorlinks=true, linkcolor=NavyBlue, urlcolor=NavyBlue, citecolor=NavyBlue]{hyperref}
\usepackage[a4paper]{geometry}
\usepackage{ts-typesetting}
\usepackage{ts-fonts}
\usepackage{booktabs}
\usepackage{csquotes}
\usepackage{float}
\usepackage{tabularx}
\usepackage[normalem]{ulem}
\usepackage{xcolor}
\usepackage{epigraph}
\usepackage{marginnote}
\usepackage{multicol}
\usepackage{progressbar}
\usepackage{graphicx}
\usepackage{tcolorbox}
\usepackage{fvextra}
\usepackage{tikz}
\usepackage{enumitem}
\usepackage{amssymb}
\usepackage{pifont}
\usepackage{amsmath}
\usepackage{seqsplit}
\newcolumntype{Y}{>{\centering\arraybackslash}X}
\newcolumntype{Z}{>{\raggedleft\arraybackslash}X}
\renewcommand*{\marginfont}{
\footnotesize
\sloppy
\emergencystretch=1em
\hyphenpenalty=50
\exhyphenpenalty=50
}
\newlength{\tsmarginparavail}
\newlength{\tsmarginparalt}
\newlength{\tsmarginparbuf}
\setlength{\tsmarginparbuf}{6mm}
\makeatletter
\AtBeginDocument{
\setlength{\tsmarginparavail}{\dimexpr\paperwidth-\textwidth-1in-\oddsidemargin-\marginparsep-\tsmarginparbuf\relax}
\ifdim\tsmarginparavail<\z@\setlength{\tsmarginparavail}{\z@}\fi
\if@twoside
\setlength{\tsmarginparalt}{\dimexpr1in+\evensidemargin-\marginparsep-\tsmarginparbuf\relax}
\ifdim\tsmarginparalt<\z@\setlength{\tsmarginparalt}{\z@}\fi
\ifdim\tsmarginparalt<\tsmarginparavail
\setlength{\tsmarginparavail}{\tsmarginparalt}
\fi
\fi
\ifdim\tsmarginparavail<\marginparwidth
\setlength{\marginparwidth}{\tsmarginparavail}
\fi
}
\makeatother
\usepackage{ts-keystrokes}
\usepackage{ts-callouts}
\usepackage{ts-code}
\usepackage{ts-glossary}
\usepackage{ts-todolist}

\usepackage{lastpage}
\usepackage{refcount}
\makeatletter
\def\ps@plain{
\let\@oddhead\@empty
\let\@evenhead\@empty
\def\@oddfoot{
\hfil
\ifnum\getpagerefnumber{LastPage}>1\relax
\thepage
\fi
\hfil}
\let\@evenfoot\@oddfoot
}
\makeatother

\title{TeXSmith Markdown Syntax}
\author{Yves Chevallier}\date{}
\begin{document}
\maketitle
\begin{abstract}
This document provides a comprehensive overview of Markdown syntax features and extensions supported by TeXSmith. It serves as a reference guide for users looking to leverage Markdown's capabilities in their documents. This is meant to be a cheatsheet of supported features and a test suite for custom templates.
\end{abstract}
\tableofcontents
\section{Core Markdown (Standard \enquote{Vanilla})}
\subsection{Headings}
Demonstrates the Markdown heading hierarchy from level 1 to level 5. Note that in Markdown only six levels are defined: \tscodeinline{\#} to \tscodeinline{\#\#\#\#\#\#}. In \LaTeX{}, however, there are less levels (e.g., \tscodeinline{\textbackslash{}section}, \tscodeinline{\textbackslash{}subsection}, \tscodeinline{\textbackslash{}subsubsection}, \tscodeinline{\textbackslash{}paragraph}, \tscodeinline{\textbackslash{}subparagraph}). The last one is usually simply converted to bold text.
\begin{tscode}[lang=md, engine=pygments]
\begin{Verbatim}[commandchars=\\\{\},breaklines, breakanywhere, commandchars=\\\{\}]
\PY{g+gh}{\PYZsh{} Heading 1}

Lorem ipsum dolor sit amet.

\PY{g+gu}{\PYZsh{}\PYZsh{} Heading 2}

Lorem ipsum dolor sit amet.

\PY{g+gu}{\PYZsh{}\PYZsh{}\PYZsh{} Heading 3}

Lorem ipsum dolor sit amet.

\PY{g+gu}{\PYZsh{}\PYZsh{}\PYZsh{}\PYZsh{} Heading 4}

Lorem ipsum dolor sit amet.

\PY{g+gu}{\PYZsh{}\PYZsh{}\PYZsh{}\PYZsh{}\PYZsh{} Heading 5}

Lorem ipsum dolor sit amet.
\end{Verbatim}
\end{tscode}
\subsection{Bold}
Shows how to emphasise text with bold weight. Simply translated to \tscodeinline{\textbackslash{}textbf\{…\}} in \LaTeX{}.
\begin{tscode}[lang=md, engine=pygments]
\begin{Verbatim}[commandchars=\\\{\},breaklines, breakanywhere, commandchars=\\\{\}]
\PY{g+gs}{**text**} or \PY{g+gs}{**t**}ex**t**
\end{Verbatim}
\end{tscode}
\begin{displayquote}
\textbf{text} or \textbf{t}ex\textbf{t}
\end{displayquote}
\subsection{Italic}
Applies emphasis for terminology or voice.
\begin{tscode}[lang=md, engine=pygments]
\begin{Verbatim}[commandchars=\\\{\},breaklines, breakanywhere, commandchars=\\\{\}]
This \PY{g+ge}{\PYZus{}text\PYZus{}} or \PY{g+ge}{*text*} is italic.
\end{Verbatim}
\end{tscode}
\begin{displayquote}
This \emph{text} or \emph{text} is italic.
\end{displayquote}
\subsection{Strikethrough}
Marks text as deleted or obsolete.
\begin{tscode}[lang=md, engine=pygments]
\begin{Verbatim}[commandchars=\\\{\},breaklines, breakanywhere, commandchars=\\\{\}]
We \PY{g+gd}{\PYZti{}\PYZti{}do not\PYZti{}\PYZti{}} want this.
\end{Verbatim}
\end{tscode}
\begin{displayquote}
We \sout{do not} want this.
\end{displayquote}
\subsection{Underline}
The \tscodeinline{\{underline\}} role underlines text. PyMdownX's \tscodeinline{\^{}\^{}text\^{}\^{}} spells the same
node, but only with the \tscodeinline{inline.insert} feature enabled; without it the carets
are literal.
\begin{tscode}[lang=md, engine=pygments]
\begin{Verbatim}[commandchars=\\\{\},breaklines, breakanywhere, commandchars=\\\{\}]
\PYZob{}underline\PYZcb{}[text]
\end{Verbatim}
\end{tscode}
\begin{displayquote}
\uline{text}
\end{displayquote}
\subsection{Inline Code / Code Blocks}
\begin{tscode}[lang=md, engine=pygments]
\begin{Verbatim}[commandchars=\\\{\},breaklines, breakanywhere, commandchars=\\\{\}]
This entry is \PY{l+s+sb}{`inline code`}, or can be in a fenced code block:

\PY{l+s+sb}{```}\PY{l+s+sb}{python}
\PY{n+nb}{print}\PY{p}{(}\PY{l+s+s2}{\PYZdq{}}\PY{l+s+s2}{Hello}\PY{l+s+s2}{\PYZdq{}}\PY{p}{)}
\PY{l+s+sb}{```}
\end{Verbatim}
\end{tscode}
\begin{displayquote}
This entry is \tscodeinline{inline code}, or can be in a fenced code block:
\begin{tscode}[lang=python, engine=pygments]
\begin{Verbatim}[commandchars=\\\{\},breaklines, breakanywhere, commandchars=\\\{\}]
\PY{n+nb}{print}\PY{p}{(}\PY{l+s+s2}{\PYZdq{}}\PY{l+s+s2}{Hello}\PY{l+s+s2}{\PYZdq{}}\PY{p}{)}
\end{Verbatim}
\end{tscode}
\end{displayquote}
\subsection{Hyperlinks}
Creates external hyperlinks, in PDF converted to clickable links.
\begin{tscode}[lang=md, engine=pygments]
\begin{Verbatim}[commandchars=\\\{\},breaklines, breakanywhere, commandchars=\\\{\}]
The German\PYZhy{}bord pysicist [\PY{n+nt}{Albert Einstein}](\PY{n+na}{https://en.wikipedia.org/wiki/Albert\PYZus{}Einstein}) is famous for
his nobel prize and the theory of relativity.
\end{Verbatim}
\end{tscode}
\begin{displayquote}
The German-bord pysicist \href{https://en.wikipedia.org/wiki/Albert\_Einstein}{Albert Einstein} is famous for
his nobel prize and the theory of relativity.
\end{displayquote}
\subsection{Images}
Embeds remote or local images.
\begin{tscode}[lang=md, engine=pygments]
\begin{Verbatim}[commandchars=\\\{\},breaklines, breakanywhere, commandchars=\\\{\}]
![\PY{n+nt}{Random Picture}](\PY{n+na}{https://picsum.photos/500/200})
\end{Verbatim}
\end{tscode}
\begin{displayquote}
\begin{figure}[H]
\centering
\includegraphics[width=\linewidth]{200.jpg}
\caption{Random Picture}
\end{figure}
\end{displayquote}
\subsection{Lists}
\subsection{Bulleted Lists}
\begin{tscode}[lang=md, engine=pygments]
\begin{Verbatim}[commandchars=\\\{\},breaklines, breakanywhere, commandchars=\\\{\}]
\PY{k}{\PYZhy{}}\PY{+w}{ }First
\PY{k}{\PYZhy{}}\PY{+w}{ }Second
\PY{+w}{    }\PY{k}{\PYZhy{}}\PY{+w}{ }Subitem
\end{Verbatim}
\end{tscode}
\begin{displayquote}
\begin{itemize}
\item First
\item Second
\begin{itemize}
\item Subitem
\end{itemize}
\end{itemize}
\end{displayquote}
\subsection{Enumerated Lists}
\begin{tscode}[lang=md, engine=pygments]
\begin{Verbatim}[commandchars=\\\{\},breaklines, breakanywhere, commandchars=\\\{\}]
\PY{k}{1.} Numbered first
\PY{k}{2.} Numbered second
\PY{+w}{    }\PY{k}{1.} Sub\PYZhy{}numbered
        a. Sub\PYZhy{}sub\PYZhy{}numbered
        b. Sub\PYZhy{}sub\PYZhy{}numbered
\PY{+w}{    }\PY{k}{2.} Sub\PYZhy{}numbered
\end{Verbatim}
\end{tscode}
\begin{displayquote}
\begin{enumerate}
\item Numbered first
\item Numbered second
\begin{enumerate}
\item Sub-numbered
a. Sub-sub-numbered
b. Sub-sub-numbered
\item Sub-numbered
\end{enumerate}
\end{enumerate}
\end{displayquote}
\subsection{Blockquotes}
Highlights quotations or cited text.
\begin{tscode}[lang=md, engine=pygments]
\begin{Verbatim}[commandchars=\\\{\},breaklines, breakanywhere, commandchars=\\\{\}]
\PY{k}{\PYZgt{} }\PY{g+ge}{Quote}
\end{Verbatim}
\end{tscode}
\begin{displayquote}
Quote
\end{displayquote}
\subsection{Horizontal Rules}
Adds a visual separator between sections.
\begin{tscode}[lang=md, engine=pygments]
\begin{Verbatim}[commandchars=\\\{\},breaklines, breakanywhere, commandchars=\\\{\}]
\PYZhy{}\PYZhy{}\PYZhy{}
\end{Verbatim}
\end{tscode}
\begin{displayquote}
\tsdivider
\end{displayquote}
\section{TeXSmith Markdown Extensions}
TeXSmith features several custom Markdown extensions to enhance document authoring.
Below is a summary of the supported extensions along with their usage.
\subsection{Small Capitals}
With standard Markdown \tscodeinline{**} or \tscodeinline{\_\_} have the same effect as bold while the former is usually preferred. TeXSmith uses the latter for small capitals when the \tscodeinline{texsmith.smallcaps}
extension is enabled.
\begin{tscode}[lang=md, engine=pygments]
\begin{Verbatim}[commandchars=\\\{\},breaklines, breakanywhere, commandchars=\\\{\}]
This is \PY{g+gs}{\PYZus{}\PYZus{}small capitals\PYZus{}\PYZus{}}.
\end{Verbatim}
\end{tscode}
\begin{displayquote}
This is \textsc{small capitals}.
\end{displayquote}
\section{Advanced Syntax Extensions (Non-Standard)}
\subsection{Tables}
\begin{tscode}[lang=md, engine=pygments]
\begin{Verbatim}[commandchars=\\\{\},breaklines, breakanywhere, commandchars=\\\{\}]
| Column 1 | Column 2 |
| \PYZhy{}\PYZhy{}\PYZhy{}\PYZhy{}\PYZhy{}\PYZhy{}\PYZhy{}\PYZhy{} | \PYZhy{}\PYZhy{}\PYZhy{}\PYZhy{}\PYZhy{}\PYZhy{}\PYZhy{}\PYZhy{} |
| Value 1  | Value 2  |
| Value 3  | Value 4  |
\end{Verbatim}
\end{tscode}
\begin{displayquote}
\begin{center}
\begin{tabularx}{\linewidth}{>{\raggedright\arraybackslash}X>{\raggedright\arraybackslash}X}
\toprule
\textbf{Column 1} & \textbf{Column 2} \\
\midrule
Value 1 & Value 2 \\
Value 3 & Value 4 \\
\bottomrule
\end{tabularx}
\end{center}
\end{displayquote}
\subsection{Footnotes}
A footnote is a note placed at the bottom of the page.
\begin{tscode}[lang=md, engine=pygments]
\begin{Verbatim}[commandchars=\\\{\},breaklines, breakanywhere, commandchars=\\\{\}]
Text with note[\PYZca{}1].

[\PY{n+nl}{\PYZca{}1}]: \PY{n+na}{Here is the note.}
\end{Verbatim}
\end{tscode}
\begin{displayquote}
Text with note\footnote{Here is the note.}.
\end{displayquote}
\subsection{Abbreviations}
\emph{Extension: \tscodeinline{abbr}}
\emph{Package: \tscodeinline{markdown} (standard)}
\emph{Description: Expands acronyms on hover.}
\begin{tscode}[lang=md, engine=pygments]
\begin{Verbatim}[commandchars=\\\{\},breaklines, breakanywhere, commandchars=\\\{\}]
The HTML standard.

*[HTML]: HyperText Markup Language
\end{Verbatim}
\end{tscode}
\begin{displayquote}
The \tsacr{HTML} standard.
\end{displayquote}
\subsection{Definition Lists}
\emph{Extension: \tscodeinline{definition\_lists}}
\emph{Package: \tscodeinline{markdown} (standard)}
\emph{Description: Pairs terms with their definitions.}
\begin{tscode}[lang=md, engine=pygments]
\begin{Verbatim}[commandchars=\\\{\},breaklines, breakanywhere, commandchars=\\\{\}]
Term
: Definition of the term

Another term
:   Another definition with a longer description that spans
    multiple lines and
    is indented properly.
\end{Verbatim}
\end{tscode}
\begin{displayquote}
\begin{description}
\item[{Term}] Definition of the term
\item[{Another term}] Another definition with a longer description that spans
multiple lines and
is indented properly.
\end{description}
\end{displayquote}
\subsection{Admonitions (Alert / Info / Tip Blocks)}
\emph{Extension: \tscodeinline{admonition}}
\emph{Package: \tscodeinline{pip install markdown} or \tscodeinline{pymdown-\allowbreak{}extensions} (for advanced styles)}
\emph{Description: Emphasises callouts such as notes and warnings.}
\begin{tscode}[lang=md, engine=pygments]
\begin{Verbatim}[commandchars=\\\{\},breaklines, breakanywhere, commandchars=\\\{\}]
::: note
This is a note.
:::

::: warning \PYZob{}collapsed=true\PYZcb{}
This is a warning.
:::
\end{Verbatim}
\end{tscode}
\tscodeinline{collapsed=true} creates a collapsible block in \tsacr{HTML} output only. PyMdownX's
\tscodeinline{!!! note} / \tscodeinline{??? warning} spell the same containers.
\begin{tscallout}[kind=note]
This is a note.
\end{tscallout}
\begin{tscallout}[kind=warning, collapsed]
This is a warning.
\end{tscallout}
\subsection{SuperFences (Enhanced Fenced Code + Nested Blocks)}
\emph{Extension: \tscodeinline{pymdownx.superfences}}
\emph{Package: \tscodeinline{pip install pymdown-\allowbreak{}extensions}}
\emph{Description: Allows nested fences (e.g., Mermaid diagrams inside fences).}
\begin{tscode}[lang=md, engine=pygments]
\begin{Verbatim}[commandchars=\\\{\},breaklines, breakanywhere, commandchars=\\\{\}]
\PY{l+s+sb}{```}\PY{l+s+sb}{mermaid}\PY{+w}{ }image width=\PYZdq{}20\PYZpc{}\PYZdq{}
\PY{l+s}{graph TD;}
\PY{l+s}{  A\PYZhy{}\PYZhy{}\PYZgt{}B;}
\PY{l+s+sb}{```}
\end{Verbatim}
\end{tscode}
\begin{displayquote}
\begin{figure}[H]
\centering
\includegraphics[width=0.2\linewidth]{<HASH>.pdf}
\end{figure}
\end{displayquote}
\subsection{Raw \LaTeX{} Blocks and Inline Snippets}
\emph{Extension: \tscodeinline{texsmith.markdown\_extensions.latex\_raw}}
\emph{Package: shipped with TeXSmith}
\emph{Description: Embeds raw \LaTeX{} that is injected verbatim.}
\begin{tscode}[lang=md, engine=pygments]
\begin{Verbatim}[commandchars=\\\{\},breaklines, breakanywhere, commandchars=\\\{\}]
\PY{l+s+sb}{```}\PY{l+s+sb}{latex}\PY{+w}{ }raw
\PY{k}{\PYZbs{}clearpage}
\PY{l+s+sb}{```}
\end{Verbatim}
\end{tscode}
\clearpage
Inline variant:
\begin{tscode}[lang=md, engine=pygments]
\begin{Verbatim}[commandchars=\\\{\},breaklines, breakanywhere, commandchars=\\\{\}]
Insert... \PYZob{}raw latex\PYZcb{}(\PYZbs{}clearpage) anywhere in the paragraph.
\end{Verbatim}
\end{tscode}
Insert... \clearpage anywhere in the paragraph.
\subsection{Emoji}
\emph{Extension: \tscodeinline{pymdownx.emoji}}
\emph{Package: \tscodeinline{pip install pymdown-\allowbreak{}extensions}}
\emph{Description: Converts shortcodes into emoji.}
\begin{tscode}[lang=md, engine=pygments]
\begin{Verbatim}[commandchars=\\\{\},breaklines, breakanywhere, commandchars=\\\{\}]
:smile:
:heart:
\end{Verbatim}
\end{tscode}
\begin{displayquote}
\tsemoji{😄}
\tsemoji{❤️}
\end{displayquote}
You can also use Unicode emoji directly (not currently supported in mono environments). It will use OpenMoji black by default in TeXSmith unless you specify another font for emoji in the configuration.
\begin{tscode}[lang=md, engine=pygments]
\begin{Verbatim}[commandchars=\\\{\},breaklines, breakanywhere, commandchars=\\\{\}]
😊 🚀 🍕 🎉 🐍 🌍 💻
📚 🎨 👽 👋 🤖 🦄 🧠
\end{Verbatim}
\end{tscode}
\begin{displayquote}
\tsemoji{😊} \tsemoji{🚀} \tsemoji{🍕} \tsemoji{🎉} \tsemoji{🐍} \tsemoji{🌍} \tsemoji{💻}
\tsemoji{📚} \tsemoji{🎨} \tsemoji{👽} \tsemoji{👋} \tsemoji{🤖} \tsemoji{🦄} \tsemoji{🧠}
\end{displayquote}
\subsection{Task Lists}
\emph{Extension: \tscodeinline{pymdownx.tasklist}}
\emph{Package: \tscodeinline{pip install pymdown-\allowbreak{}extensions}}
\emph{Description: Creates interactive checklists.}
\begin{tscode}[lang=md, engine=pygments]
\begin{Verbatim}[commandchars=\\\{\},breaklines, breakanywhere, commandchars=\\\{\}]
\PY{k}{\PYZhy{} }\PY{k}{[x]} Done
\PY{k}{\PYZhy{} }\PY{k}{[ ]} To do
\end{Verbatim}
\end{tscode}
\begin{displayquote}
\begin{tstasklist}
\item[\tsdone] Done
\item[\tstodo] To do
\end{tstasklist}
\end{displayquote}
\subsection{Highlight / Mark}
\emph{Extension: \tscodeinline{pymdownx.mark}}
\emph{Package: \tscodeinline{pip install pymdown-\allowbreak{}extensions}}
\emph{Description: Highlights text with a marker effect.}
\begin{tscode}[lang=md, engine=pygments]
\begin{Verbatim}[commandchars=\\\{\},breaklines, breakanywhere, commandchars=\\\{\}]
==Cellular respiration is a set of metabolic reactions== that take place
in the cells of organisms. Its ==primary== function is to convert
biochemical energy from nutrients into ==adenosine triphosphate (ATP)==,
and then release waste products.
\end{Verbatim}
\end{tscode}
\begin{displayquote}
\tsmark{Cellular respiration is a set of metabolic reactions} that take place in the
cells of organisms. Its \tsmark{primary} function is to convert biochemical energy
from nutrients into \tsmark{adenosine triphosphate (ATP)}, and then release waste
products.
\end{displayquote}
\subsection{Tilde / Subscript / Superscript}
\emph{Extension: \tscodeinline{pymdownx.tilde}, \tscodeinline{pymdownx.caret}}
\emph{Package: \tscodeinline{pip install pymdown-\allowbreak{}extensions}}
\emph{Description: Adds subscripts and superscripts.}
\begin{tscode}[lang=md, engine=pygments]
\begin{Verbatim}[commandchars=\\\{\},breaklines, breakanywhere, commandchars=\\\{\}]
H\PYZti{}2\PYZti{}O, E = mc\PYZca{}2\PYZca{}

H₂O, E = mc²
\end{Verbatim}
\end{tscode}
\begin{displayquote}
H\textsubscript{2}O, E = mc\textsuperscript{2}

H\textsubscript{2}O, E = mc\textsuperscript{2}
\end{displayquote}
\subsection{SmartSymbols}
\emph{Extension: \tscodeinline{pymdownx.smartsymbols}}
\emph{Package: \tscodeinline{pip install pymdown-\allowbreak{}extensions}}
\emph{Description: Automatically substitutes typographic punctuation.}
\begin{tscode}[lang=md, engine=pygments]
\begin{Verbatim}[commandchars=\\\{\},breaklines, breakanywhere, commandchars=\\\{\}]
\PYZdq{}Smart quotes\PYZdq{}, ellipses..., en\PYZhy{}dash \PYZhy{}\PYZhy{}, em\PYZhy{}dash \PYZhy{}\PYZhy{}\PYZhy{}
\end{Verbatim}
\end{tscode}
\begin{displayquote}
\enquote{Smart quotes}, ellipses..., en-dash --, em-dash ---
\end{displayquote}
\subsection{Better Math / \LaTeX{}}
\emph{Extension: \tscodeinline{pymdownx.arithmatex}}
\emph{Package: \tscodeinline{pip install pymdown-\allowbreak{}extensions}}
\emph{Rendering engines: depends on KaTeX or MathJax.}
\emph{Description: Typesets mathematical formulas with \LaTeX{} syntax.}
\begin{tscode}[lang=md, engine=pygments]
\begin{Verbatim}[commandchars=\\\{\},breaklines, breakanywhere, commandchars=\\\{\}]
The differential form of Maxwell\PYZsq{}s equations:

\PYZdl{}\PYZdl{}
\PYZbs{}begin\PYZob{}aligned\PYZcb{}
\PYZbs{}nabla \PYZbs{}cdot \PYZbs{}mathbf\PYZob{}E\PYZcb{} \PYZam{}= \PYZbs{}frac\PYZob{}\PYZbs{}rho\PYZcb{}\PYZob{}\PYZbs{}varepsilon\PYZus{}0\PYZcb{} \PYZbs{}\PYZbs{}
\PYZbs{}nabla \PYZbs{}cdot \PYZbs{}mathbf\PYZob{}B\PYZcb{} \PYZam{}= 0 \PYZbs{}\PYZbs{}
\PYZbs{}nabla \PYZbs{}times \PYZbs{}mathbf\PYZob{}E\PYZcb{} \PYZam{}= \PYZhy{}\PYZbs{}frac\PYZob{}\PYZbs{}partial \PYZbs{}mathbf\PYZob{}B\PYZcb{}\PYZcb{}\PYZob{}\PYZbs{}partial t\PYZcb{} \PYZbs{}\PYZbs{}
\PYZbs{}nabla \PYZbs{}times \PYZbs{}mathbf\PYZob{}B\PYZcb{} \PYZam{}= \PYZbs{}mu\PY{g+ge}{\PYZus{}0 \PYZbs{}mathbf\PYZob{}J\PYZcb{} + \PYZbs{}mu\PYZus{}}0 \PYZbs{}varepsilon\PYZus{}0
\PYZbs{}frac\PYZob{}\PYZbs{}partial \PYZbs{}mathbf\PYZob{}E\PYZcb{}\PYZcb{}\PYZob{}\PYZbs{}partial t\PYZcb{}
\PYZbs{}end\PYZob{}aligned\PYZcb{}
\PYZdl{}\PYZdl{}
\end{Verbatim}
\end{tscode}
\begin{displayquote}
The differential form of Maxwell's equations:

$$
\begin{aligned}
\nabla \cdot \mathbf{E} &= \frac{\rho}{\varepsilon_0} \\
\nabla \cdot \mathbf{B} &= 0 \\
\nabla \times \mathbf{E} &= -\frac{\partial \mathbf{B}}{\partial t} \\
\nabla \times \mathbf{B} &= \mu_0 \mathbf{J} + \mu_0 \varepsilon_0
\frac{\partial \mathbf{E}}{\partial t}
\end{aligned}
$$
\end{displayquote}
\subsection{MagicLink (Automatic GitHub / Issue Links)}
\emph{Extension: \tscodeinline{pymdownx.magiclink}}
\emph{Package: \tscodeinline{pip install pymdown-\allowbreak{}extensions}}
\emph{Description: Auto-links URLs, emails, GitHub references, and more.}
\begin{tscode}[lang=md, engine=pygments]
\begin{Verbatim}[commandchars=\\\{\},breaklines, breakanywhere, commandchars=\\\{\}]
https://github.comuser/project\PYZsh{}1
\end{Verbatim}
\end{tscode}
\begin{displayquote}
\url{https://github.comuser/project\#1}
\end{displayquote}
\subsection{ProgressBar}
\emph{Extension: \tscodeinline{pymdownx.progressbar}}
\emph{Package: \tscodeinline{pip install pymdown-\allowbreak{}extensions}}
\emph{Description: Visualises progress as textual bars.}
\begin{tscode}[lang=markdown, engine=pygments]
\begin{Verbatim}[commandchars=\\\{\},breaklines, breakanywhere, commandchars=\\\{\}]
[=25\PYZpc{} \PYZdq{}Research\PYZdq{}]
[=50\PYZpc{} \PYZdq{}Implementation\PYZdq{}]
[=75\PYZpc{} \PYZdq{}Review\PYZdq{}]
[=100\PYZpc{} \PYZdq{}Launch\PYZdq{}]\PYZob{}.thin\PYZcb{}
\end{Verbatim}
\end{tscode}
\tsprogress{0.25}{Research}
\tsprogress{0.5}{Implementation}
\tsprogress{0.75}{Review}
\tsprogress[thin]{1}{Launch}
\subsection{Details / Collapsible Blocks}
\emph{Extension: \tscodeinline{pymdownx.details}}
\emph{Package: \tscodeinline{pip install pymdown-\allowbreak{}extensions}}
\emph{Description: Creates collapsible disclosure sections.}
\begin{tscode}[lang=md, engine=pygments]
\begin{Verbatim}[commandchars=\\\{\},breaklines, breakanywhere, commandchars=\\\{\}]
::: note \PYZob{}title=Title collapsed=false\PYZcb{}
Collapsible content.
:::
\end{Verbatim}
\end{tscode}
\begin{displayquote}
\begin{tscallout}[kind=note, title={Title}]
Collapsible content.
\end{tscallout}
\end{displayquote}
\subsection{Keys (Keyboard Display)}
\emph{Extension: \tscodeinline{pymdownx.keys}}
\emph{Package: \tscodeinline{pip install pymdown-\allowbreak{}extensions}}
\emph{Description: Shows keyboard shortcuts with consistent styling.}
\begin{tscode}[lang=md, engine=pygments]
\begin{Verbatim}[commandchars=\\\{\},breaklines, breakanywhere, commandchars=\\\{\}]
++Ctrl+C++
\end{Verbatim}
\end{tscode}
\begin{displayquote}
\tskeys{Ctrl,C}
\end{displayquote}
\subsection{Tab Blocks (Tabbed Content)}
\emph{Extension: \tscodeinline{pymdownx.tabbed}}
\emph{Package: \tscodeinline{pip install pymdown-\allowbreak{}extensions}}
\emph{Description: Groups content in tabbed panes.}
\begin{tscode}[lang=md, engine=pygments]
\begin{Verbatim}[commandchars=\\\{\},breaklines, breakanywhere, commandchars=\\\{\}]
:::: tabs
::: tab \PYZob{}title=Windows\PYZcb{}
Windows is a Microsoft operating system.
:::
::: tab \PYZob{}title=Linux\PYZcb{}
Linux is an open\PYZhy{}source operating system.
:::
::::
\end{Verbatim}
\end{tscode}
\begin{tsdiv}{tab}[title=Windows]
Windows is a Microsoft operating system.
\end{tsdiv}
\begin{tsdiv}{tab}[title=Linux]
Linux is an open-source operating system.
\end{tsdiv}
\subsection{Meta-Data / Front Matter}
\emph{Extension: \tscodeinline{meta}}
\emph{Package: \tscodeinline{markdown} (standard)}
\emph{Description: Adds YAML metadata at the beginning of a document.}
\begin{tscode}[lang=md, engine=pygments]
\begin{Verbatim}[commandchars=\\\{\},breaklines, breakanywhere, commandchars=\\\{\}]
\PYZhy{}\PYZhy{}\PYZhy{}
Title: My document
\PY{g+gu}{Author: Alice}
\PY{g+gu}{\PYZhy{}\PYZhy{}\PYZhy{}}

\PY{g+gh}{\PYZsh{} Title}
\end{Verbatim}
\end{tscode}
\subsection{Snippets}
\emph{Extension: \tscodeinline{pymdownx.snippets}}
\emph{Package: \tscodeinline{pymdown-\allowbreak{}extensions}}
\emph{Description: Includes content from other Markdown files.}

Permits you to include content from other Markdown files.
\begin{tscode}[lang=md, engine=pygments]
\begin{Verbatim}[commandchars=\\\{\},breaklines, breakanywhere, commandchars=\\\{\}]
\PY{l+s+sb}{```}\PY{l+s+sb}{python}\PY{+w}{ }include=\PYZdq{}hanoi.py\PYZdq{}
\PY{l+s+sb}{```}
\end{Verbatim}
\end{tscode}
\begin{tscode}[lang=python, engine=pygments]
\begin{Verbatim}[commandchars=\\\{\},breaklines, breakanywhere, commandchars=\\\{\}]
\PY{k+kn}{from}\PY{+w}{ }\PY{n+nn}{\PYZus{}\PYZus{}future\PYZus{}\PYZus{}}\PY{+w}{ }\PY{k+kn}{import} \PY{n}{annotations}

\PY{k+kn}{from}\PY{+w}{ }\PY{n+nn}{collections}\PY{n+nn}{.}\PY{n+nn}{abc}\PY{+w}{ }\PY{k+kn}{import} \PY{n}{Iterable}

\PY{n}{Move} \PY{o}{=} \PY{n+nb}{tuple}\PY{p}{[}\PY{n+nb}{int}\PY{p}{,} \PY{n+nb}{str}\PY{p}{,} \PY{n+nb}{str}\PY{p}{]}

\PY{k}{def}\PY{+w}{ }\PY{n+nf}{tower\PYZus{}of\PYZus{}hanoi}\PY{p}{(}\PY{n}{n}\PY{p}{:} \PY{n+nb}{int}\PY{p}{,} \PY{n}{source}\PY{p}{:} \PY{n+nb}{str}\PY{p}{,} \PY{n}{destination}\PY{p}{:} \PY{n+nb}{str}\PY{p}{,} \PY{n}{auxiliary}\PY{p}{:} \PY{n+nb}{str}\PY{p}{)} \PY{o}{\PYZhy{}}\PY{o}{\PYZgt{}} \PY{n+nb}{list}\PY{p}{[}\PY{n}{Move}\PY{p}{]}\PY{p}{:}
\PY{+w}{    }\PY{l+s+sd}{\PYZdq{}\PYZdq{}\PYZdq{}Compute the move sequence for the Tower of Hanoi puzzle.\PYZdq{}\PYZdq{}\PYZdq{}}
    \PY{n}{moves}\PY{p}{:} \PY{n+nb}{list}\PY{p}{[}\PY{n}{Move}\PY{p}{]} \PY{o}{=} \PY{p}{[}\PY{p}{]}

    \PY{k}{def}\PY{+w}{ }\PY{n+nf}{\PYZus{}solve}\PY{p}{(}\PY{n}{disks}\PY{p}{:} \PY{n+nb}{int}\PY{p}{,} \PY{n}{start}\PY{p}{:} \PY{n+nb}{str}\PY{p}{,} \PY{n}{end}\PY{p}{:} \PY{n+nb}{str}\PY{p}{,} \PY{n}{spare}\PY{p}{:} \PY{n+nb}{str}\PY{p}{)} \PY{o}{\PYZhy{}}\PY{o}{\PYZgt{}} \PY{k+kc}{None}\PY{p}{:}
        \PY{k}{if} \PY{n}{disks} \PY{o}{==} \PY{l+m+mi}{0}\PY{p}{:}
            \PY{k}{return}
        \PY{n}{\PYZus{}solve}\PY{p}{(}\PY{n}{disks} \PY{o}{\PYZhy{}} \PY{l+m+mi}{1}\PY{p}{,} \PY{n}{start}\PY{p}{,} \PY{n}{spare}\PY{p}{,} \PY{n}{end}\PY{p}{)}
        \PY{n}{moves}\PY{o}{.}\PY{n}{append}\PY{p}{(}\PY{p}{(}\PY{n}{disks}\PY{p}{,} \PY{n}{start}\PY{p}{,} \PY{n}{end}\PY{p}{)}\PY{p}{)}
        \PY{n}{\PYZus{}solve}\PY{p}{(}\PY{n}{disks} \PY{o}{\PYZhy{}} \PY{l+m+mi}{1}\PY{p}{,} \PY{n}{spare}\PY{p}{,} \PY{n}{end}\PY{p}{,} \PY{n}{start}\PY{p}{)}

    \PY{n}{\PYZus{}solve}\PY{p}{(}\PY{n}{n}\PY{p}{,} \PY{n}{source}\PY{p}{,} \PY{n}{destination}\PY{p}{,} \PY{n}{auxiliary}\PY{p}{)}
    \PY{k}{return} \PY{n}{moves}

\PY{k}{def}\PY{+w}{ }\PY{n+nf}{render\PYZus{}solution}\PY{p}{(}\PY{n}{moves}\PY{p}{:} \PY{n}{Iterable}\PY{p}{[}\PY{n}{Move}\PY{p}{]}\PY{p}{)} \PY{o}{\PYZhy{}}\PY{o}{\PYZgt{}} \PY{n+nb}{str}\PY{p}{:}
\PY{+w}{    }\PY{l+s+sd}{\PYZdq{}\PYZdq{}\PYZdq{}Return a human\PYZhy{}readable description of the move sequence.\PYZdq{}\PYZdq{}\PYZdq{}}
    \PY{k}{return} \PY{l+s+s2}{\PYZdq{}}\PY{l+s+se}{\PYZbs{}n}\PY{l+s+s2}{\PYZdq{}}\PY{o}{.}\PY{n}{join}\PY{p}{(}
        \PY{l+s+sa}{f}\PY{l+s+s2}{\PYZdq{}}\PY{l+s+s2}{Move disk }\PY{l+s+si}{\PYZob{}}\PY{n}{disk}\PY{l+s+si}{\PYZcb{}}\PY{l+s+s2}{ from source }\PY{l+s+si}{\PYZob{}}\PY{n}{start}\PY{l+s+si}{\PYZcb{}}\PY{l+s+s2}{ to destination }\PY{l+s+si}{\PYZob{}}\PY{n}{end}\PY{l+s+si}{\PYZcb{}}\PY{l+s+s2}{\PYZdq{}} \PY{k}{for} \PY{n}{disk}\PY{p}{,} \PY{n}{start}\PY{p}{,} \PY{n}{end} \PY{o+ow}{in} \PY{n}{moves}
    \PY{p}{)}
\end{Verbatim}
\end{tscode}
\subsection{EscapeAll}
\emph{Extension: \tscodeinline{pymdownx.escapeall}}
\emph{Package: \tscodeinline{pymdown-\allowbreak{}extensions}}
\emph{Description: Forces every Markdown character to be escaped.}

Forces the escaping of special Markdown characters.
\subsection{BetterEm}
\emph{Extension: \tscodeinline{pymdownx.betterem}}
\emph{Package: \tscodeinline{pymdown-\allowbreak{}extensions}}
*Description: Manages combinations of bold and italic (\tscodeinline{***text***}, etc.).*
\begin{tscode}[lang=md, engine=pygments]
\begin{Verbatim}[commandchars=\\\{\},breaklines, breakanywhere, commandchars=\\\{\}]
\PY{g+gs}{**\PYZus{}bold and italic\PYZus{}**}

***bold and italic***
\end{Verbatim}
\end{tscode}
\begin{displayquote}
\tslead{\emph{bold and italic}}

\emph{\textbf{bold and italic}}
\end{displayquote}
\subsection{Long Dash (---)}
\emph{Description: Notes how em-dashes are handled differently by parsers.}

No current extension automatically handles the em dash in Markdown. In \LaTeX{} you use \tscodeinline{-\allowbreak{}-\allowbreak{}} for an en dash and \tscodeinline{-\allowbreak{}-\allowbreak{}-\allowbreak{}} for an em dash. The \tscodeinline{-\allowbreak{}-\allowbreak{}-\allowbreak{}} sequence can be confusing in Markdown because it is also used for horizontal rules; however, some parsers such as \tscodeinline{markdown-\allowbreak{}it-\allowbreak{}py} interpret it contextually. Experiment as needed.
\begin{tscode}[lang=md, engine=pygments]
\begin{Verbatim}[commandchars=\\\{\},breaklines, breakanywhere, commandchars=\\\{\}]
Achiles \PYZhy{}\PYZhy{} the swifest runner \PYZhy{}\PYZhy{} was fast, but not as fast as the tortoise.
\end{Verbatim}
\end{tscode}
\begin{displayquote}
Achiles -- the swifest runner -- was fast, but not as fast as the tortoise.
\end{displayquote}
\subsection{Additional Syntax}
\emph{Description: Highlights extra syntaxes such as callouts and directives.}
\begin{tscode}[lang=md, engine=pygments]
\begin{Verbatim}[commandchars=\\\{\},breaklines, breakanywhere, commandchars=\\\{\}]
\PY{k}{\PYZgt{} }\PY{g+ge}{[!note] This is a note.}
\PY{k}{\PYZgt{} }\PY{g+ge}{Used on Docusaurus, Obsidian, GitHub.}

::directive\PYZob{}param\PYZcb{} Used on MDX, Astro
\end{Verbatim}
\end{tscode}
\subsection{\textbf{Other, Less Common Dialects}}
\emph{Description: Summarises less common Markdown dialects and their strengths.}
\begin{center}
\begin{tabularx}{\linewidth}{>{\raggedright\arraybackslash}X>{\raggedright\arraybackslash}X>{\raggedright\arraybackslash}X}
\toprule
\textbf{Dialect} & \textbf{Highlights} & \textbf{Package} \\
\midrule
\textbf{MultiMarkdown} & Tables, footnotes, citations, extended metadata & \tscodeinline{multimarkdown} \\
\textbf{RMarkdown} & Executable code (knitr), math, tables & \tscodeinline{rmarkdown} (R) \\
\textbf{kramdown} & Math, footnotes, block IAL, definition lists & (Ruby) \\
\textbf{CommonMark Ext.} & Full GFM support & \tscodeinline{cmarkgfm} \\
\textbf{Ghost / Jekyll} & ``Liquid tags'' (\tscodeinline{\{\% include \%\}}) & Liquid engine \\
\textbf{MkDocs Material} & PyMdownX support + Snippets + Tabs + Admonitions & \tscodeinline{mkdocs-\allowbreak{}material} \\
\bottomrule
\end{tabularx}
\end{center}
\printglossary[type=\acronymtype]
\end{document}
